Salah satu kemampuan yang harus dimiliki oleh seorang Sarjana Teknik adalah kemampuan menganalisis sumber-sumber pustaka. Oleh karena itu, di dalam dokuman C-251 ini, mahasiswa perlu melakukan studi pustaka dan menuliskannya sampai kepada level kedalaman dan detail yang memadai. Sebelum melakukan analisis pustaka, uraikan lagi apa saja permasalahan yang ingin Anda selesaikan di dalam \textit{capstone} ini. Dari permasalahan tersebut, carilah literatur yang sudah pernah membahas dan menyelesaikan masalah yang sama, dan pilihkan pustaka kunci yang akan Anda adopsi untuk membantu menyelesaikan Anda. 

Untuk melakukan analisis studi, tentukan minimal tiga pustaka yang paling berkaitan dengan solusi yang diusulkan. Buatlah analisis yang mendalam dan jika perlu disimulasikan metodenya untuk menyelesaikan permasalahan Anda. Uraikan secara ringkas sumber-sumber yang membahas berbagai kerangka teori dan metode-metode yang mungkin untuk digunakan secara sistematis. Gunakan \textit{heading} untuk menguraikan metode yang ada.

\section{Metode 1}
\label{sec:Metode_01}

    \lipsum[1-3]

\section{Metode 2}
\label{sec:Metode_02}

    \lipsum[4-6]
    
    \subsection{Sub-Metode 2}
    \label{subsec:Sub-Metode_02}
    
    \lipsum[7-9]

\section{Metode 3}
\label{sec:Metode_03}

    \lipsum[10-12]
    
    \subsection{Sub-Metode 3A}
    \label{subsec:Sub-Metode_03A}
    
    \lipsum[13-15]
    
    \subsection{Sub-Metode 3B}
    \label{subsec:Sub-Metode_03B}
    
    \lipsum[15-18]


\section{Pemilihan Metode}
\label{sec:Pemilihan_Metode}
    
    Setelah dideskripsikan kekurangan dan kelebihan metode, tim akan memilih metode yang paling tepat disertai pengembangan metode tersebut. Jika dari semua metode yang ada tidak menghasilkan solusi yang konkret, tim bisa membuat metode sendiri. Dalam hal ini, proyek \textit{capstone} ini akan menjalankan mode penelitian sehingga tidak diwajibkan untuk menghasilkan luaran tertentu. Walaupun begitu adanya luaran seperti tertulis pada lembar luaran sangat disarankan. 
    
    \textbf{\uline{Apabila metode yang diimplementasikan berbeda dengan apa yang Anda usulkan dalam dokumen C-251}}, maka berikan penjelasan detail mengenai perubahan metode tersebut. Sampaikanlah pula argumen/justifikasi mengapa perubahan metode tersebut perlu Anda lakukan dalam proses implementasi.